\section{What the Qualitative Evidence Suggests About Mechanisms}
\label{sec:mechanisms}

We discuss mechanisms cautiously as hypotheses constrained by audits, not as proven causal pathways.
The goal is to connect RQ5--RQ6 to the quantitative tables without over-fitting a single story to every configuration.

\paragraph{Hypothesis A: lexical/register calibration.}
Priming may increase the density of corpus-consistent ISV lexical choices and reduce Polish-residue forms, moving canonical coverage without rewriting source structure.
LOW and UNSEEN positives with retained plot/exposition are compatible with this hypothesis.
Under this view, the corpus functions less as a narrative template and more as a register/lexicon prior exposed in context.

\paragraph{Hypothesis B: orthographic normalization.}
Outside-inventory reductions often co-occur with positive $\Delta$, suggesting that orthographic cleanup is one pathway by which the metric can move.
Divergent cases show it is neither necessary nor sufficient.
Metric--ortho divergence is itself scientifically useful: it warns against treating any single diagnostic as a complete account of Primed behaviour.

\paragraph{Hypothesis C: corpus-opening reuse under high overlap.}
On HIGH, large positives recurrently co-occur with winter-opening near-copy in the qualitative audit.
This may inflate coverage when corpus lexis enters the output, but it is not the only route to positive $\Delta$, and opening copy can appear with negative $\Delta$ (Qwen HIGH r03).
Hypothesis~C therefore explains a HIGH-local signature rather than the whole experiment.

\paragraph{Compatibility, not competition.}
A, B, and C are not mutually exclusive.
A configuration can show orthographic cleanup and lexical borrowing; a HIGH run can show opening reuse plus lexical adoption.
The important empirical constraint is that A/B-compatible behaviour on LOW/UNSEEN indicates that opening-copy (Hypothesis~C) is not necessary for observing a positive coverage $\Delta$.

\paragraph{What we do not claim.}
We do not claim that models ``learned Interslavic'' from the corpus in a demonstrated causal sense, that priming is the sole cause of $\Delta$, or that metric gains equal communicative quality.
A useful descriptive framing---not a formal law---is that priming appears to be a \emph{context-sensitive intervention} whose observed effect depends on the interaction between model configuration and source regime.

\paragraph{Consequences for future experiments.}
If Hypothesis~C dominates in a future high-overlap narrative arm, experiments should measure opening overlap explicitly rather than treating it as anecdotal.
If Hypotheses~A/B dominate in low-overlap or unseen arms, evaluation should track lexical inventory hits and orthography separately from narrative similarity.
Phase~2B already separates those diagnostics; future designs can pre-register which diagnostic is primary for which regime.
